\section{FORMAL RESTATEMENT}
\textbf{Erd\H{o}s \#31; reconstruction, 2026-09-23.}
Given an infinite set $A\subseteq\mathbb N$, construct a set
$B\subseteq\mathbb N$ with natural density zero such that $A+B$
contains every sufficiently large integer. We use positive integers
throughout; adjoining zero to either set does not affect the conclusion.

\section{QUICK LITERATURE AND CONTEXT CHECK}
The current page credits Lorentz with proving the Erd\H{o}s--Straus
conjecture and labels the result proved, with a formal verification
reported. The following finite-cover and interval-patching argument
reconstructs the existence result. No novelty or independent verification
of the reported formal proof is claimed.

\section{WORK}
\subsection{1. A finite cyclic covering lemma}
Let $F\subseteq\mathbb Z/q\mathbb Z$ have $j\geq1$ elements. There
exists $R\subseteq\mathbb Z/q\mathbb Z$ with $F+R=\mathbb Z/q\mathbb Z$
and
\[
 |R|\leq q\eta_j,\qquad \eta_j=\frac{\log j+1}{j}.    \tag{1}
\]
For $j=1$ take all residues. For $j\geq2$, choose a random set $R_0$
by including each residue independently with probability $p=(\log j)/j$.
Each target residue is missed by $F+R_0$ with probability
$(1-p)^j\leq e^{-pj}=1/j$. Fix $f_0\in F$. For each missed target $x$,
add $x-f_0$ to $R_0$. The repaired set covers every target and has
expected size at most $qp+q/j=q\eta_j$. Some outcome satisfies (1).
This is a finite existence argument; exhaustive search would select
such a set without randomness if desired.

\subsection{2. Covers with densities tending to zero}
List $A$ increasingly as $a_1<a_2<\cdots$, set $a_0=a_1$, and let
$F_j=\{a_1,\ldots,a_j\}$, $m_j=a_j$, and $q_j=m_j+1$.
The elements of $F_j$ are distinct modulo $q_j$. Choose $R_j$ by (1),
and let $P_j$ be the periodic set of integers with residues in $R_j$.
Then every integer has a representation $a+b$ with $a\in F_j$ and
$b\in P_j$, and $P_j$ has density at most $\eta_j\to0$.
For any interval of $L$ consecutive integers,
\[
 |P_j\cap I|\leq\eta_j L+q_j,                         \tag{2}
\]
by separating complete periods and a final incomplete period.

\subsection{3. Patching the periodic covers}
Put $H_j=q_j+m_j+a_0+1$. Choose increasing positive integer thresholds
$T_j$ so large that
\[
 T_j\geq j^2\sum_{i=1}^jH_i,\qquad T_1>a_0,
 \qquad T_{j+1}>2(T_j+m_j+a_0+1).                     \tag{3}
\]
This is possible recursively. Define
\[
 B=\bigcup_{j\geq1}\left(
       P_j\cap[T_j,T_{j+1})\right)
   \ \cup\ \bigcup_{j\geq1}
       \bigl([T_j-a_0,T_j+m_j]\cap\mathbb Z\bigr).    \tag{4}
\]
All elements are positive by the choice of thresholds. The second union
repairs the left boundary of each interval.

Indeed, given $n\geq T_1$, choose $j$ with $T_j\leq n<T_{j+1}$.
If $n<T_j+m_j$, then $n=a_0+(n-a_0)$ and the second summand belongs
to the boundary interval in (4). Otherwise choose $a\in F_j$ with
$n-a\in P_j$. Since $a\leq m_j$, we have
$T_j\leq n-a<T_{j+1}$, so this second summand belongs to the first
union in (4). Thus every $n\geq T_1$ lies in $A+B$.

\subsection{4. Density at every cutoff, including interval boundaries}
Let $T_J\leq x<T_{J+1}$. By (2), the contribution of the first union
to $B\cap[1,x]$ is at most
\[
 \sum_{j\leq J}\eta_j L_j(x)+\sum_{j\leq J}q_j,
\]
where $L_j(x)=|[T_j,T_{j+1})\cap[1,x]\cap\mathbb Z|$ and
$\sum_jL_j(x)\leq x$. The boundary intervals with $j\leq J$ add at
most $\sum_{j\leq J}(m_j+a_0+1)$. The interval for $J+1$, if it has
already begun, adds at most $a_0$ integers below $T_{J+1}$; later
boundary intervals contribute nothing, by (3). Consequently
\[
 |B\cap[1,x]|\leq\sum_{j\leq J}\eta_j L_j(x)
                   +\sum_{j\leq J}H_j+a_0
 \leq\sum_{j\leq J}\eta_j L_j(x)+\frac{x}{J^2}+a_0.
\]
Given $\epsilon>0$, choose $j_0$ so that $\eta_j<\epsilon$ for
$j\geq j_0$. The finitely many earlier intervals contribute a fixed
constant $C_{j_0}$, while the remaining sum is at most $\epsilon x$.
It follows that
\[
 \limsup_{x\to\infty}\frac{|B\cap[1,x]|}{x}\leq\epsilon.
\]
Since $\epsilon$ was arbitrary, $B$ has natural density zero.
\hfill$\square$

\section{VERIFICATION AND LIMITS}
The proof requires neither positive density nor a growth bound for $A$.
Large gaps in $A$ are handled by choosing larger thresholds in (3).
The error terms include both incomplete periods and boundary repairs;
showing small density only at the thresholds would not have sufficed.
The finite covering lemma is checked exhaustively on cyclic groups
of order at most eight in the accompanying batch script.

\section{FINAL STATUS}
\textbf{SOLVED (known affirmative answer reconstructed).}
The construction and the density estimate are complete. They establish
existence and a finite-search description relative to an enumeration
of $A$, without claiming an efficient algorithm or new quantitative
bounds for the complement.

\section{SOURCES}
T.~F.~Bloom, problem \#31, with attribution to G.~G.~Lorentz (1954),
\url{https://www.erdosproblems.com/31}, checked 2026-09-23.
Checks: \texttt{verification/solved\_batch\_26\_checks.py}.

This estimates completion of the stated existence proof, not novelty.

\section{COMPLETION ESTIMATE}
\textbf{COMPLETION: 100\%}
